% Tema: partida de motor / hardware (aula 04)
\question[3] Por que um botão ligado a uma entrada digital que não tem resistor de pull interno precisa de um resistor de pull-down (ou pull-up) externo?

\begin{solution}
Sem um resistor de pull, quando o botão está solto a entrada fica ``flutuando'' (não conectada firmemente nem ao nível alto nem ao baixo), captando ruído eletromagnético do ambiente e podendo ler valores aleatórios e instáveis. O resistor de pull-down (ligado ao GND) ou pull-up (ligado à alimentação) garante um nível lógico definido quando o botão está solto, e o botão, ao ser pressionado, sobrepõe esse nível com o oposto.
\end{solution}
